Una empresa minera extrae \textbf{80 toneladas} de mineral rojo y \textbf{60 toneladas} de mineral negro cada semana. Estos pueden ser tratados de diferentes maneras para producir tres diferentes aleaciones, blanda, dura o fuerte. Para producir 1 tonelada de aleación blanda se requiere \textbf{4 toneladas} de mineral rojo y \textbf{2 toneladas} de negro. Para la aleación dura los requisitos son \textbf{1 toneladas} de rojo y \textbf{4 toneladas} de negro, mientras que para la aleación fuerte son \textbf{2 toneladas} de rojo y \textbf{4 toneladas} de negro. La ganancia por tonelada de la venta de las aleaciones (después de considerar los costos de producción, pero no los de extracción, que se consideran fijos) son \textbf{300}, \textbf{200} y \textbf{300 euros} por tonelada de aleación blanda, dura y fuerte, respectivamente.

Además, el mineral rojo se puede vender a \textbf{50} euros la tonelada y el mineral negro se debe emplear al completo.

Existe un compromiso comercial de servir al menos \textbf{10} toneladas conjuntamente entre las aleaciones dura y fuerte.

\begin{equation}
\begin{split}
\mbox{max. } z = 300x_1 + 200x_2 + 300x_3 + 50x_4\\
s.a.:\\
4x_1 + 1x_2 + 2x_3 + 1x_4 \leq 80\\
2x_1 + 4x_2 + 4x_3 + 0x_4 = 60\\
0x_1 + 1x_2 + 1x_3 + 0x_4 \geq 10\\
x_1,\,\,x_2,\,\,x_3,\,\,x_4\geq 0\\
\end{split}
\end{equation}

Donde $x_1$, $x_2$ y $x_3$ representan las producciones de aleación blanda, dura y fuerte (en toneladas), respectivamente y $x_4$ son las toneladas de mineral rojo de venta en el mercado.

La siguiente tabla es la correspondiente a la solución óptima del problema original

\begin{center}
\begin{tabular}{c|c|cccccccc|}
 & $z$ & $x_1$ & $x_2$ & $x_3$ & $x_4$ & $h_1$ & $h_3$ & $a_2$ & $a_3$\\ 
Phase 2 & -7000 & 0 & -50 & 0 & 0 & -50 & 0 & -50 & 0 \\ 
\hline
$x_4$ & 20  & 0 & -1 & 0 & 1 & 1 & -6 & -2 & 6\\ 
$x_1$ & 10  & 1 & 0 & 0 & 0 & 0 & 2 & 1/2 & -2\\ 
$x_3$ & 10  & 0 & 1 & 1 & 0 & 0 & -1 & 0 & 1\\ 
\hline
\end{tabular}
\end{center}


\begin{enumerate}

\item \textbf{Interpretar} el resultado y explicar el plan de extracción y el significado de los valores de las variables de holgura.


\item Dado que el mineral negro se debe emplear completo, \textbf{justificar} si la empresa estaría interesada en extraer más o menos cantidad que las 60 toneladas semanales.

\item \textbf{Justificar haciendo uso de los fundamentos del método del simplex} por qué, sin necesidad de resolver el problema, se podría afirmar que en la solución óptima $h_1$ nunca puede tomar un valor diferente de 0.

\item Cuál es el intervalo relativo a la cantidad del compromiso comercial (aleaciones dura y fuerte) dentro del cual se mantiene la gama de producción correspondiente a la solución óptima.

\item La empresa está interesada en saber cuál sería el impacto sobre su plan de producción y sobre su beneficio si, una determinada semana, no pudiera vender (y, por lo tanto, no produciría) más de 8 toneladas de aleación fuerte.

% \begin{comment}
% Podemos preguntar sobre el V^B d ela variable artificiales o su signo o su relación con la restricción...
% \end{comment}

\end{enumerate}


% \begin{comment}
% \begin{center}
% \begin{tabular}{c|c|cccccc|}
%  & $z$ & $x_1$ & $x_2$ & $h_1$ & $h_3$ & $a_2$ & $a_3$\\ 
% Phase 1 & 0 & 0 & 0 & 0 & 0 & 0 & -1 \\ 
% Phase 2 & -45 & 0 & -1/4 & 0 & 0 & -3/4 & 0 \\ 
% \hline
% $h_1$ & 50  & 0 & 5/2 & 1 & 0 & -1/2 & 0\\ 
% $h_3$ & 5  & 0 & -1/4 & 0 & 1 & 1/4 & -1\\ 
% $x_1$ & 15  & 1 & 3/4 & 0 & 0 & 1/4 & 0\\ 
% \hline
% \end{tabular}
% \end{center}

% \end{comment}